\begin{problem}[第30届全国中学生物理竞赛复赛][运动介质中的电磁场]{5}
\begin{wrapfigure}{r}{0.5\textwidth}
    \vspace*{-1.0\baselineskip} 
    \centering
    \includegraphics[width=0.9\linewidth]{image/1.png}
    \vspace*{-5.5\baselineskip}   % 在底部向上收拢
\end{wrapfigure}
平行板电容器两极板分别位于 $z = \pm \frac{d}{2}$ 的平面内，电容器起初未被充电。整个装置处于均匀磁场中，磁感应强度大小为 $B$，方向沿 $x$ 轴负方向。
\begin{subproblem}
在电容器参考系 $S$ 中只存在磁场；而在以沿 $y$ 轴正方向的恒定速度 $(0, v, 0)$（这里 $(0, v, 0)$ 表示为沿 $x$、$y$、$z$ 轴正方向的速度分量分别为 $0$、$v$、$0$，以下类似）相对于电容器运动的参考系 $S'$ 中，可能既有电场 $(E_x', E_y', E_z')$ 又有磁场 $(B_x', B_y', B_z')$。试在非相对论情形下，从伽利略速度变换，求出在参考系 $S'$ 中电场 $(E_x', E_y', E_z')$ 和磁场 $(B_x', B_y', B_z')$ 的表达式。已知电荷量和作用在物体上的合力在伽利略变换下不变。
\end{subproblem}

\begin{subproblem}
现在让介电常数为 $\varepsilon$ 的电中性液体（绝缘体）在平行板电容器两极板之间匀速流动，流速大小为 $v$，方向沿 $y$ 轴正方向。在相对液体静止的参考系（即相对于电容器运动的参考系）$S'$ 中，由于液体处在第1问所述的电场 $(E_x', E_y', E_z')$ 中，其正负电荷会因电场力作用而发生相对移动（即所谓极化效应），使得液体中出现附加的静电感应电场，因而液体中总电场强度不再是 $(E_x', E_y', E_z')$，而是 $\frac{\varepsilon_0}{\varepsilon}(E_x', E_y', E_z')$，这里 $\varepsilon_0$ 是真空的介电常数。这将导致在电容器参考系 $S$ 中电场不再为零。试求电容器参考系 $S$ 中电场的强度以及电容器上、下极板之间的电势差。（结果用 $\varepsilon_0$、$\varepsilon$、$v$、$B$ 或（和）$d$ 表出。）
\end{subproblem}
\end{problem}